% Section 7: the bosonic Fock space inside the formal Fock space, coefficient-finite maps and their
% maximal realizations, closed forms (Kato VI 1.3), product vectors and the photon-added class.
% GA-BMPS milestone B4.5, general part (20 Oct 2026). New text for vetting (vetbenv).
\providecommand{\dom}{\operatorname{dom}}
\providecommand{\form}{\mathfrak q}
\begin{vetbenv}
\section{The bosonic Fock space and closed forms}\label{sec:fockspace}

The bosonic Fock space is the space of square-summable vectors of the formal Fock space of
Section~\ref{sec:fock}. A linear map whose coefficients are finite linear combinations of
coefficients acts on every formal vector, and its restriction to the vectors of the Fock space
that it maps back into the Fock space is a closed operator. Finite sums of squared norms of such
maps are therefore closed quadratic forms in the sense of Kato~\cite[VI\,\S1.3]{kato}, the setting
of the representation theorem for self-adjoint operators~\cite[VI\,\S2]{kato}; \vetc{the parent
forms of the bosonic matrix product states of~\cite{tjoacirac} are of this kind, and, where they are densely defined, their operators are the parent
Hamiltonians}.

\medskip
\noindent\textbf{Notation for this section.}
\begin{center}
\begin{tabular}{@{}L{0.3\textwidth}L{0.64\textwidth}@{}}
$\N^{(\iota)}$ & the multi-indices $n=(n_i)_{i\in\iota}$, finitely many nonzero (Section~\ref{sec:fock})\\
$\lVert\psi\rVert^2$ & squared norm of a formal vector, in $[0,\infty]$ (Definition~\ref{def:fock-hilbert})\\
$\Gamma_s(\iota)$, $\langle\varphi,\psi\rangle$, $\jmath$ & bosonic Fock space, its inner product, its
embedding into $\Fock(\iota,\C)$ (Definition~\ref{def:fock-hilbert})\\
$\mathcal S_n$ & a finite set of multi-indices (Definition~\ref{def:fock-coefffinite})\\
$\bigotimes_j\chi_j$, $\psi_{j,n}$, $S_j$ & product vector, mode-$j$ fiber, a single-mode map $S$ acting
on mode $j$ (Definition~\ref{def:fock-modes})\\
$\bar T$, $\dom\bar T$ & maximal realization of $T$ and its domain (Definition~\ref{def:fock-maxreal})\\
$(T_k)_{k\in\mathcal K}$, $\dom\form_\Gamma$, $\form_\Gamma$ & a finite family of maps, its form domain and
its form (Definition~\ref{def:fock-form})\\
$\ket{q,\kappa,\ell}$, $\mathcal P_{\kappa,\ell}$ & photon-added Gaussian state and photon-added class
(Definition~\ref{def:fock-photon})\\
\end{tabular}
\end{center}
In this section $\iota$ is any set of modes unless it is said to be finite, $T$ and $T'$ are linear
maps on $\Fock(\iota,\C)$, $S$ is a linear map on the single-mode space $\Fock(\{1\},\C)$, whose
multi-indices we identify with $\N$, and $p$ is a polynomial, not the exponent $p\ge1$ of
Section~\ref{sec:fock}.

\subsection{The bosonic Fock space}

\begin{definition}[Squared norm; bosonic Fock space]\label{def:fock-hilbert}
Let $\iota$ be a set of modes.
\begin{enumerate}
\item[(a)] For $\psi\in\Fock(\iota,\C)$ the \emph{squared norm} is $\lVert\psi\rVert^2=\sum_n\lvert
\psi(n)\rvert^2\in[0,\infty]$, the sum of the non-negative terms over all multi-indices.
\item[(b)] The \emph{bosonic Fock space} $\Gamma_s(\iota)=\ell^2(\N^{(\iota)})$ is the Hilbert space of
the square-summable families $(\psi(n))_n$ over the set $\N^{(\iota)}$ of multi-indices $n$, with the inner
product $\langle\varphi,\psi\rangle=\sum_n\overline{\varphi(n)}\,\psi(n)$, antilinear in $\varphi$. The
\emph{embedding} $\jmath:\Gamma_s(\iota)\to\Fock(\iota,\C)$ is $\jmath(\psi)(n)=\psi(n)$.
\end{enumerate}
\end{definition}

The number states $\ket n$ form an orthonormal basis of $\Gamma_s(\iota)$. For finite $\iota$ this is
the space $\Gamma_s(\mathfrak h)$, $\mathfrak h=\C^\iota$, of Section~\ref{sec:fock} written in the
number basis; for infinite $\iota$ it is in the same way the Fock space over $\ell^2(\iota)$~\cite{dg}.
The formal vector with the coefficients $\sqrt{(2k)!}/k!$ at $\ket{2k}$ of Section~\ref{sec:fock} has
infinite squared norm.

\begin{lemma}\label{lem:fock-hilbert}
Let $\iota$ be a set of modes.
\begin{enumerate}
\item[(a)] $\jmath$ is linear and injective, and $\lVert\jmath\psi\rVert^2=\lVert\psi\rVert^2$ for
$\psi\in\Gamma_s(\iota)$, the right side being the squared Hilbert norm.
\item[(b)] A formal vector $\varphi\in\Fock(\iota,\C)$ lies in $\jmath(\Gamma_s(\iota))$ if and only if
$\lVert\varphi\rVert^2<\infty$.
\item[(c)] The finitely supported vectors, the finite linear combinations of number states, are dense
in $\Gamma_s(\iota)$.
\end{enumerate}
\end{lemma}

\begin{proof}
\begin{enumerate}
\item[(a)] Linearity holds coefficientwise, and $\jmath\psi$ has the coefficients of $\psi$. The squared
Hilbert norm $\sum_n\lvert\psi(n)\rvert^2$ is the real sum of a summable non-negative family, which
equals its sum in $[0,\infty]$.
\item[(b)] A sum of non-negative terms in $[0,\infty]$ is finite exactly when the family of terms is
summable, that is, when $\varphi$ is square-summable.
\item[(c)] For $\psi\in\Gamma_s(\iota)$ and a finite set $F$ of multi-indices,
$\lVert\psi-\sum_{n\in F}\psi(n)\ket n\rVert^2=\sum_{n\notin F}\lvert\psi(n)\rvert^2$, which is as small
as we like for $F$ large, by summability.\qedhere
\end{enumerate}
\end{proof}

\subsection{Coefficient-finite maps and maximal realizations}\label{sec:fock-coefffinite}

\begin{definition}[Coefficient-finite]\label{def:fock-coefffinite}
A linear map $T$ on $\Fock(\iota,\C)$ is \emph{coefficient-finite} if for every multi-index $n$ there
is a finite set $\mathcal S_n$ of multi-indices such that $(T\psi)(n)=(T\varphi)(n)$ whenever
$\psi(m)=\varphi(m)$ for every $m\in\mathcal S_n$.
\end{definition}

Since $\psi$ agrees on $\mathcal S_n$ with $\sum_{m\in\mathcal S_n}\psi(m)\ket m$, linearity gives
$(T\psi)(n)=\sum_{m\in\mathcal S_n}T_{nm}\,\psi(m)$ with $T_{nm}=(T\ket m)(n)$. A coefficient-finite map
is thus one whose matrix in the number basis is \emph{row-finite}, and every row-finite matrix
defines a coefficient-finite map by this formula. The ladder operators are of this kind, and so is
every exponential of Definition~\ref{def:fock-raising} on finitely many modes.

\begin{definition}[Product vectors; operators acting on one mode]\label{def:fock-modes}
Let $\iota$ be a set of modes.
\begin{enumerate}
\item[(a)] For finite $\iota$ and $\chi_j\in\Fock(\{1\},\C)$, $j\in\iota$, the \emph{product vector}
$\bigotimes_j\chi_j\in\Fock(\iota,\C)$ has the coefficients $\bigl(\bigotimes_j\chi_j\bigr)(n)=\prod_j
\chi_j(n_j)$.
\item[(b)] For $\psi\in\Fock(\iota,\C)$, $j\in\iota$ and a multi-index $n$, the \emph{mode-$j$ fiber}
$\psi_{j,n}\in\Fock(\{1\},\C)$ is $\psi_{j,n}(m)=\psi\bigl(n+(m-n_j)e_j\bigr)$, the coefficient at $n$ with
$n_j$ replaced by $m$.
\item[(c)] For a linear map $S$ on $\Fock(\{1\},\C)$ and $j\in\iota$, the \emph{map $S$ acting on mode
$j$} is the linear map $S_j$ on $\Fock(\iota,\C)$ with $(S_j\psi)(n)=(S\psi_{j,n})(n_j)$.
\end{enumerate}
\end{definition}

For example, $a$ acting on mode $j$ is $a_j$: $(a\psi_{j,n})(n_j)=\sqrt{n_j+1}\,\psi(n+e_j)$. The fiber
does not depend on $n_j$, and the fiber of a product vector along $j$ is a multiple of its factor at
$j$.

\Needspace{10\baselineskip}
\begin{lemma}[Coefficient-finite maps]\label{lem:fock-coefffinite}
Let $\iota$ be a set of modes.
\begin{enumerate}
\item[(a)] The maps $0$, $1$, $a_j$ and $a_j^\dagger$ are coefficient-finite. If $T$ and $T'$ are
coefficient-finite, so are $T+T'$, $TT'$, $cT$ for $c\in\C$, every power $T^k$, every finite sum of
coefficient-finite maps and $p(T)$ for every polynomial $p\in\C[X]$. If $S$ is a coefficient-finite
map on $\Fock(\{1\},\C)$, then $S_j$ is coefficient-finite for every $j\in\iota$. If $T$ is
coefficient-finite and raises the particle number, then $e^T$ is coefficient-finite.
\item[(b)] If $\iota$ is finite, every linear map on $\Fock(\iota,\C)$ that raises the particle number
is coefficient-finite.
\item[(c)] If $T$ is coefficient-finite, then for every multi-index $n$ the map $\psi\mapsto
(T\jmath\psi)(n)$ is continuous on $\Gamma_s(\iota)$.
\end{enumerate}
\end{lemma}

\begin{proof}
\begin{enumerate}
\item[(a)] For $0$ and $1$ take $\mathcal S_n=\emptyset$ and $\{n\}$; for $a_j$ take $\{n+e_j\}$, and for
$a_j^\dagger$ take $\{n-e_j\}$ if $n_j\ge1$ and $\emptyset$ otherwise (Definition~\ref{def:fock-ladder}).
For $T+T'$ and $cT$ the union of the two sets serves. For $TT'$, $(TT'\psi)(n)$ depends on the
coefficients $(T'\psi)(m)$, $m\in\mathcal S^T_n$, each of which depends on $\psi$ on
$\mathcal S^{T'}_m$; the union of these finitely many finite sets serves. Powers follow by induction,
and finite sums and $p(T)$ by the first claims. For $S_j$, $(S_j\psi)(n)$ depends on $\psi_{j,n}$ at the
finitely many $m\in\mathcal S^S_{n_j}$, that is, on $\psi$ at the multi-indices $n+(m-n_j)e_j$. For
$e^T$, $(e^T\psi)(n)=\sum_{r\le\lvert n\rvert}(T^r\psi)(n)/r!$ (Definition~\ref{def:fock-raising}) depends
on $\psi$ on the union of the sets of the maps $T^r$, $r\le\lvert n\rvert$.
\item[(b)] By Lemma~\ref{lem:fock-raising}(b), $(T\psi)(n)$ depends only on the coefficients $\psi(n')$
with $\lvert n'\rvert<\lvert n\rvert$, and for finite $\iota$ there are finitely many such $n'$. In
particular every polynomial in the creation operators with zero constant term is coefficient-finite
(Lemma~\ref{lem:fock-raising}(c)).
\item[(c)] $(T\jmath\psi)(n)=\sum_{m\in\mathcal S_n}T_{nm}\,\psi(m)$ is a finite linear combination of
coordinates, and $\lvert\psi(m)-\psi'(m)\rvert\le\lVert\psi-\psi'\rVert$ in $\Gamma_s(\iota)$.\qedhere
\end{enumerate}
\end{proof}

An operator $A$ on $\Gamma_s(\iota)$, defined on a subspace $\dom A$, is \emph{closed} if its graph
$\{(\psi,A\psi):\psi\in\dom A\}$ is closed in $\Gamma_s(\iota)\times\Gamma_s(\iota)$, and \emph{densely
defined} if $\dom A$ is dense~\cite[III\,\S5]{kato}.

\begin{definition}[Maximal realization]\label{def:fock-maxreal}
For a linear map $T$ on $\Fock(\iota,\C)$, the \emph{maximal realization} $\bar T$ is the operator on
$\Gamma_s(\iota)$ with the domain $\dom\bar T=\{\psi\in\Gamma_s(\iota):\lVert T\jmath\psi\rVert^2<\infty\}$
and $\bar T\psi$ the vector with $\jmath(\bar T\psi)=T\jmath\psi$ (Lemma~\ref{lem:fock-hilbert}(b)).
\end{definition}

The domain is a subspace, since a sum of square-summable vectors is square-summable. For a
row-finite matrix, $\bar T$ is the maximal operator of the matrix on $\ell^2$: the largest operator on
$\Gamma_s(\iota)$ that acts on each vector of its domain by the matrix. No domain has to be chosen.

\begin{theorem}[Closedness of the maximal realization]\label{thm:fock-maxreal-closed}
Let $T$ be coefficient-finite. Then $\bar T$ is closed: if $\psi^{(r)}\in\dom\bar T$, $\psi^{(r)}\to
\psi$ and $\bar T\psi^{(r)}\to\chi$ in $\Gamma_s(\iota)$, then $\psi\in\dom\bar T$ and $\bar T\psi=\chi$.
\end{theorem}

\begin{proof}
Fix $n$. By Lemma~\ref{lem:fock-coefffinite}(c), $(T\jmath\psi^{(r)})(n)\to(T\jmath\psi)(n)$, and since
coordinates are continuous on $\Gamma_s(\iota)$, $(T\jmath\psi^{(r)})(n)=(\bar T\psi^{(r)})(n)\to\chi(n)$.
Hence $T\jmath\psi=\jmath\chi$, which has finite squared norm; so $\psi\in\dom\bar T$ and $\bar T\psi=\chi$.
\end{proof}

\subsection{Closed forms}\label{sec:fock-closedforms}

We use the forms of~\cite[VI\,\S1]{kato} with the arguments in the order of the inner product. A
\emph{form} $\mathfrak t$ on $\Gamma_s(\iota)$ is a map $\mathfrak t(\varphi,\psi)$ on $\dom\mathfrak t\times
\dom\mathfrak t$, for a subspace $\dom\mathfrak t$, the \emph{form domain}, antilinear in $\varphi$ and
linear in $\psi$; it is \emph{Hermitian} if $\overline{\mathfrak t(\varphi,\psi)}=\mathfrak t(\psi,\varphi)$
and \emph{non-negative} if $\mathfrak t(\psi,\psi)\ge0$. A non-negative Hermitian form is \emph{closed}
if $\psi^{(r)}\in\dom\mathfrak t$, $\psi^{(r)}\to\psi$ and $\mathfrak t(\psi^{(r)}-\psi^{(s)},\psi^{(r)}-
\psi^{(s)})\to0$ as $r,s\to\infty$ imply $\psi\in\dom\mathfrak t$ and $\mathfrak t(\psi^{(r)}-\psi,
\psi^{(r)}-\psi)\to0$~\cite[VI\,\S1.3]{kato}; equivalently, $\dom\mathfrak t$ is complete in the norm
$(\lVert\psi\rVert^2+\mathfrak t(\psi,\psi))^{1/2}$~\cite[VI, Theorem~1.11]{kato}. It is \emph{densely
defined} if $\dom\mathfrak t$ is dense, and a \emph{form core} is a subspace of $\dom\mathfrak t$ dense in
that norm.

\begin{definition}[The form of a finite family]\label{def:fock-form}
Let $(T_k)_{k\in\mathcal K}$ be a finite family of linear maps on $\Fock(\iota,\C)$. Its \emph{form
domain} is $\dom\form_\Gamma=\bigcap_k\dom\bar T_k$, and its \emph{form} is
\[
 \form_\Gamma(\varphi,\psi)=\sum_{k\in\mathcal K}\langle\bar T_k\varphi,\bar T_k\psi\rangle,\qquad
 \varphi,\psi\in\dom\form_\Gamma .
\]
\end{definition}

On the formal Fock space the same family has the \emph{formal form} $\sum_k\lVert T_k\varphi\rVert^2\in
[0,\infty]$, defined for every $\varphi\in\Fock(\iota,\C)$. The next lemma identifies $\form_\Gamma$ with
the formal form read on $\Gamma_s(\iota)$.

\begin{lemma}[Form domain and form]\label{lem:fock-form}
Let $(T_k)_{k\in\mathcal K}$ be a finite family of linear maps on $\Fock(\iota,\C)$.
\begin{enumerate}
\item[(a)] $\dom\form_\Gamma$ is a subspace of $\Gamma_s(\iota)$, and $\dom\form_\Gamma=\{\psi\in
\Gamma_s(\iota):\sum_k\lVert T_k\jmath\psi\rVert^2<\infty\}$.
\item[(b)] $\form_\Gamma$ is antilinear in $\varphi$, linear in $\psi$ and Hermitian, and on the diagonal
it is real and equals the formal form: $\form_\Gamma(\psi,\psi)=\sum_k\lVert T_k\jmath\psi\rVert^2$ for
$\psi\in\dom\form_\Gamma$. In particular $\form_\Gamma$ is non-negative.
\end{enumerate}
\end{lemma}

\begin{proof}
\begin{enumerate}
\item[(a)] An intersection of subspaces is a subspace, and a finite sum in $[0,\infty]$ is finite
exactly when each term is.
\item[(b)] Each term is the inner product composed with the linear map $\bar T_k$ in both arguments,
which gives the linearity and $\overline{\langle\bar T_k\varphi,\bar T_k\psi\rangle}=\langle\bar
T_k\psi,\bar T_k\varphi\rangle$. On the diagonal, by Lemma~\ref{lem:fock-hilbert}(a),
\[
 \langle\bar T_k\psi,\bar T_k\psi\rangle=\lVert\bar T_k\psi\rVert^2=\lVert T_k\jmath\psi\rVert^2 .\qedhere
\]
\end{enumerate}
\end{proof}

\begin{figure}[ht]
\centering\vetbcolor
\begin{tikzpicture}
 \draw[frame] (0,0) rectangle (12,5);
 \node[anchor=north west] at (0.1,4.95) {$\Fock(\iota,\C)$};
 \draw[region] (5,2.4) ellipse (4.6 and 2.1);
 \node at (5,4.05) {$\jmath(\Gamma_s(\iota))$: $\lVert\psi\rVert^2<\infty$};
 \begin{scope}
  \clip (3.8,2.1) ellipse (2.2 and 1.1);
  \fill[gray!20] (5.6,2.1) ellipse (2.2 and 1.1);
 \end{scope}
 \draw[region] (3.8,2.1) ellipse (2.2 and 1.1);
 \draw[region] (5.6,2.1) ellipse (2.2 and 1.1);
 \node at (2.4,2.25) {$\dom\bar T_1$};
 \node at (6.65,2.5) {$\dom\bar T_2$};
 \node at (4.7,2.45) {$\dom\form_\Gamma$};
 \node[dot] (p) at (4.7,1.7) {}; \node[right=2pt] at (p) {$\psi$};
 \node[dot] (tp) at (4.6,0.72) {}; \node[right=2pt] at (tp) {$T_1\psi$};
 \draw[map,shorten <=1pt,shorten >=1pt] (p) -- (tp);
 \node[dot] (q) at (7.0,1.7) {}; \node[left=2pt] at (q) {$\psi'$};
 \node[dot] (t1q) at (10.8,1.2) {}; \node[below right=1pt] at (t1q) {$T_1\psi'$};
 \draw[map,shorten <=1pt,shorten >=1pt] (q) to[bend right=10] (t1q);
 \node[dot] (t2q) at (8.05,2.65) {}; \node[right=2pt] at (t2q) {$T_2\psi'$};
 \draw[map,shorten <=1pt,shorten >=1pt] (q) -- (t2q);
\end{tikzpicture}
\caption{The form domain of a family $T_1,T_2$ (Definition~\ref{def:fock-form}). The maps act on the whole
formal Fock space; the domain of the maximal realization $\bar T_k$ consists of the vectors of
$\Gamma_s(\iota)$ that $T_k$ maps back into $\Gamma_s(\iota)$, and the form domain is the intersection of
these domains. Here $\psi'$ lies in $\dom\bar T_2$ but not in $\dom\bar T_1$: $T_2\psi'$ is square-summable and
$T_1\psi'$ is not.}\label{fig:fock-formdomain}
\end{figure}

\begin{theorem}[Closed form]\label{thm:fock-form-closed}
Let $(T_k)_{k\in\mathcal K}$ be a finite family of coefficient-finite maps on $\Fock(\iota,\C)$. If
$\psi^{(r)}\in\dom\form_\Gamma$, $\psi^{(r)}\to\psi$ in $\Gamma_s(\iota)$ and $\sum_k\lVert
T_k\jmath(\psi^{(r)}-\psi^{(s)})\rVert^2\to0$ as $r,s\to\infty$, then $\psi\in\dom\form_\Gamma$ and
$\sum_k\lVert T_k\jmath(\psi^{(r)}-\psi)\rVert^2\to0$. So $\form_\Gamma$ is closed.
\end{theorem}

\begin{proof}
Each term of the sum tends to zero, and $\lVert T_k\jmath(\psi^{(r)}-\psi^{(s)})\rVert=\lVert\bar
T_k\psi^{(r)}-\bar T_k\psi^{(s)}\rVert$ by linearity and Lemma~\ref{lem:fock-hilbert}(a), so each sequence
$\bar T_k\psi^{(r)}$ is Cauchy in $\Gamma_s(\iota)$ and converges to some $\chi_k$. By
Theorem~\ref{thm:fock-maxreal-closed}, $\psi\in\dom\bar T_k$ and $\bar T_k\psi=\chi_k$. Hence
$\psi\in\dom\form_\Gamma$, and $\lVert T_k\jmath(\psi^{(r)}-\psi)\rVert^2=\lVert\bar T_k\psi^{(r)}-\chi_k
\rVert^2\to0$ for each of the finitely many $k$. By Lemma~\ref{lem:fock-form}(b) this is closedness in
the sense above.
\end{proof}

\vetx{The theorem applies two results of Kato to the closed operators $\bar T_k$ of Theorem~\ref{thm:fock-maxreal-closed}: the form
$\lVert T\psi\rVert^2$ of a closed operator $T$, on its domain, is closed~\cite{kato}, and a finite sum
of closed non-negative forms, on the intersection of their domains, is closed~\cite{kato}; the proof
above checks this directly.} Closedness needs no hypothesis beyond coefficient-finiteness; whether $\form_\Gamma$ is densely defined
depends on the family.

\subsection{Product vectors and the photon-added class}\label{sec:fock-photon}

\begin{lemma}[Squared norm of a product vector]\label{lem:fock-tensornorm}
\begin{enumerate}
\item[(a)] For a finite set $\iota$, a set $A$ and functions $g_i:A\to[0,\infty]$, $i\in\iota$,
\[
 \sum_{f:\iota\to A}\ \prod_{i\in\iota}g_i(f(i))=\prod_{i\in\iota}\ \sum_{a\in A}g_i(a)\quad\text{in }[0,\infty],
\]
with the convention $0\cdot\infty=0$.
\item[(b)] For finite $\iota$ and $\chi_j\in\Fock(\{1\},\C)$, $\lVert\bigotimes_j\chi_j\rVert^2=\prod_j
\lVert\chi_j\rVert^2$ in $[0,\infty]$.
\end{enumerate}
\end{lemma}

\begin{proof}
\begin{enumerate}
\item[(a)] Induction on the number of elements of $\iota$. For $\iota=\emptyset$ both sides are $1$.
Adding an element $i_0$, the functions on $\iota\cup\{i_0\}$ are the pairs of a value $a$ at $i_0$ and a
function on $\iota$; sums of non-negative terms in $[0,\infty]$ may be taken in any order and are
multiplied termwise by a constant factor, so the sum is $\sum_ag_{i_0}(a)$ times the sum over $\iota$.
\item[(b)] On finitely many modes the multi-indices are the functions $\iota\to\N$, and
$\lvert\prod_j\chi_j(n_j)\rvert^2=\prod_j\lvert\chi_j(n_j)\rvert^2$; apply (a) with $A=\N$.\qedhere
\end{enumerate}
\end{proof}

For a single mode and $\kappa,\ell\in\C$, the map $\kappa(a^\dagger)^2+\ell a^\dagger$ raises the particle
number (Lemma~\ref{lem:fock-raising}(c)); its exponential $e^{\kappa(a^\dagger)^2+\ell a^\dagger}$ is the
\emph{Gaussian operator}; the name is used here only for these exponentials of creation operators, not for
Gaussian unitaries.

\begin{definition}[Photon-added Gaussian states; the photon-added class]\label{def:fock-photon}
For $q\in\N$ and $\kappa,\ell\in\C$, the \emph{photon-added Gaussian state} is $\ket{q,\kappa,\ell}=
(a^\dagger)^q\,e^{\kappa(a^\dagger)^2+\ell a^\dagger}\ket0\in\Fock(\{1\},\C)$, and the \emph{photon-added
class} is $\mathcal P_{\kappa,\ell}=\operatorname{span}\{\ket{q,\kappa,\ell}:q\in\N\}$.
\end{definition}

\vetc{The state $\ket{0,\kappa,\ell}$ is a formal vector, which for $\lvert\kappa\rvert<\tfrac12$ is a multiple of a
squeezed coherent state; the states $\ket{q,\kappa,\ell}$ are the
photon-added Gaussian states of~\cite[Sec.~II\,A\,1]{tjoacirac}\vetx{; for $\kappa=0$ they are multiples of the photon-added
coherent states of Agarwal and Tara~\cite{agarwaltara}}.} The maps $\kappa(a^\dagger)^2$ and
$\ell a^\dagger$ commute, so by Theorem~\ref{thm:fock-explaw} the Gaussian operator is
$e^{\kappa(a^\dagger)^2}e^{\ell a^\dagger}$, and by Theorems~\ref{thm:fock-exp-bargmann}
and~\ref{thm:fock-bargmann} the Bargmann series of $\ket{q,\kappa,\ell}$ is $z^qE$ with
$E=e^{\kappa z^2}e^{\ell z}$. So $\mathcal P_{\kappa,\ell}$ consists of the formal vectors whose Bargmann
series are $p(z)E$, $p$ a polynomial.

\begin{lemma}[The photon-added class]\label{lem:fock-photon}
Let $\kappa,\ell\in\C$.
\begin{enumerate}
\item[(a)] $\ket{q,\kappa,\ell}=(a^\dagger)^q\ket{0,\kappa,\ell}$, and $\mathcal P_{\kappa,\ell}=\{p(a^\dagger)
\ket{0,\kappa,\ell}:p\in\C[X]\}$.
\item[(b)] $\mathcal P_{\kappa,\ell}$ is invariant under $a$ and under $a^\dagger$, hence under $p(a)$ and
under $p(a^\dagger)$ for every polynomial $p$.
\item[(c)] The Gaussian operator $e^{\kappa(a^\dagger)^2+\ell a^\dagger}$ maps $\mathcal P_{0,0}$ into
$\mathcal P_{\kappa,\ell}$.
\item[(d)] $\ket{q,0,0}=\sqrt{q!}\,\ket q$ for every $q$, so every number state lies in $\mathcal P_{0,0}$.
\end{enumerate}
\end{lemma}

\begin{proof}
The Bargmann map is injective, so identities may be checked on Bargmann series, where $a^\dagger$ is
multiplication by $z$, $a$ is the derivative $\partial$ and $p(a^\dagger)$ is multiplication by $p$
(Theorem~\ref{thm:fock-bargmann}).
\begin{enumerate}
\item[(a)] The first claim is the definition, and a linear combination of the $(a^\dagger)^q\ket{0,\kappa,
\ell}$ is $p(a^\dagger)\ket{0,\kappa,\ell}$ for the polynomial with these coefficients.
\item[(b)] $a^\dagger p(a^\dagger)=(Xp)(a^\dagger)$. For $a$: from $e^{cz^m}=\sum_kc^kz^{mk}/k!$ the
coefficientwise derivative is $\partial e^{cz^m}=mc\,z^{m-1}e^{cz^m}$ for $m\ge1$, so the Leibniz rule gives
$\partial E=(2\kappa z+\ell)E$ and $\partial(pE)=\bigl(p'+(2\kappa z+\ell)p\bigr)E$. Hence $a\,p(a^\dagger)
\ket{0,\kappa,\ell}=\bigl(p'+(2\kappa X+\ell)p\bigr)(a^\dagger)\ket{0,\kappa,\ell}$. A subspace invariant under
a map is invariant under its powers and their linear combinations.
\item[(c)] The Bargmann series of $e^{\kappa(a^\dagger)^2+\ell a^\dagger}p(a^\dagger)\ket0$ is $E\,p$, that of
$p(a^\dagger)\ket{0,\kappa,\ell}$.
\item[(d)] The exponential of the zero map is the identity (Definition~\ref{def:fock-raising}), so
$\ket{q,0,0}=(a^\dagger)^q\ket0=\sqrt{q!}\,\ket q$ (after Theorem~\ref{thm:fock-ccr}).\qedhere
\end{enumerate}
\end{proof}

By (b) the class is also invariant under every finite sum of products of $a$ and $a^\dagger$. It
contains no information about norms: $\ket{0,\kappa,0}$ is square-summable exactly when
$\lvert\kappa\rvert<\tfrac12$, its squared moduli being $\binom{2k}{k}\lvert\kappa\rvert^{2k}$ at $\ket{2k}$.

\subsection*{Notes and references}

Closed forms, form cores and the representation theorems that turn a closed, densely defined,
non-negative form into a self-adjoint operator are those of Kato~\cite[VI\,\S\S1--2]{kato}. The forms here
are built from maximal realizations, which act on every vector they can, so no dense domain has to be
chosen and closedness reduces to the continuity of the coefficients (Lemma~\ref{lem:fock-coefffinite}(c)).
The usual construction starts from the operators on a dense domain $\mathcal D$, such as the finitely
supported vectors, and takes the closure of the form $\sum_k\lVert T_k\psi\rVert^2$ on $\mathcal D$. When
$\mathcal D\subseteq\dom\form_\Gamma$, that closure is the restriction of the closed form $\form_\Gamma$ of
Theorem~\ref{thm:fock-form-closed} to the closure of $\mathcal D$ in the norm $(\lVert\psi\rVert^2+
\form_\Gamma(\psi,\psi))^{1/2}$, and the two forms agree exactly when $\mathcal D$ is a form core.
\end{vetbenv}
